%% ============================================================
%% preamble/fonts.tex — engine-conditional fontspec, 3 profiles
%%
%% Purpose      : loads the font stack selected by book.yaml's
%%                typography.font_profile (\BookFontProfile):
%%                  libertinus — Libertinus Serif/Sans/Mono (default)
%%                  garamond   — EB Garamond + Source Sans 3 (falls
%%                               back to TeX Gyre Heros) + Noto Sans Mono
%%                  plex       — IBM Plex Serif/Sans/Mono, SemiBold-as-bold
%%                Also defines the small-caps helpers used by styling
%%                (\bkscshape, \bklabelcaps) and the cover font switches
%%                (\coverTitleFont, \coverSubtitleFont, \coverBodyFont).
%% Dependencies : etoolbox (\ifdefstring), generated/metadata.tex
%%                (\BookFontProfile). Loads fontspec itself. microtype
%%                settings are guarded so this module also works under
%%                latex/cover/cover.tex (article class, minimal preamble).
%% Provenance   : ADR 0004. libertinus stack from datacenter/wiki/
%%                the-last-book; garamond from complexity/htsd/legal-tech;
%%                plex (incl. the SemiBold-as-bold print calibration and
%%                the protrusion table) from history-through-rfc-book.md.
%%
%% RULES:
%%   * Fonts are loaded by FAMILY NAME, never by file path — hard-coded
%%     paths broke htsd's portability (ADR 0004).
%%   * LuaLaTeX resolves texmf-dist fonts by name via luaotfload;
%%     XeLaTeX needs the fonts visible to fontconfig. `make doctor`
%%     checks the active profile resolves before a full build.
%% ============================================================

\RequirePackage{iftex}
\ifPDFTeX
  \PackageError{book-template}
    {Font profiles are OpenType-only; compile with LuaLaTeX or XeLaTeX}{}
\fi
\RequirePackage{fontspec}

% TeX ligatures (---, --, ``'') everywhere, in every profile.
\defaultfontfeatures{Ligatures=TeX}

% ----------------------------------------------------------------------
% Helpers styling.tex relies on. Profiles that lack a real small-caps
% face override these (Plex ships no smcp feature — see plex branch).
%   \bkscshape      — font switch for small-caps contexts
%   \bklabelcaps{x} — even small-caps label (running heads, "chapter one")
% ----------------------------------------------------------------------
\newcommand{\bkscshape}{\scshape}
\newcommand{\bklabelcaps}[1]{{\scshape\MakeLowercase{#1}}}

% ======================================================================
% Profile: libertinus (default)
% Body/sans/mono all from the Libertinus family — a single super-family
% keeps color and x-height coherent with zero scaling fiddles.
% ======================================================================
\ifdefstring{\BookFontProfile}{libertinus}{%
  % Oldstyle (text) figures in running prose; tables switch to lining
  % below. Libertinus carries both onum and lnum/tnum.
  \setmainfont{Libertinus Serif}[Numbers=OldStyle]
  \setsansfont{Libertinus Sans}
  % Libertinus Mono ships regular only; fake the slant/weight so code
  % comments (\itshape) and keywords (\bfseries) don't fall back to
  % upright with a font-shape warning.
  \setmonofont{Libertinus Mono}[Scale=MatchLowercase,
    AutoFakeSlant=0.18, AutoFakeBold=1.5]
  % Lining + tabular figures inside every tabular: digits must align
  % in columns even though prose uses oldstyle (ADR 0004).
  \AtBeginEnvironment{tabular}{\addfontfeatures{Numbers={Monospaced,Lining}}}
  % Cover identity (consumed by latex/cover/cover.tex).
  \newcommand{\coverTitleFont}{\rmfamily\bfseries}
  \newcommand{\coverSubtitleFont}{\sffamily}
  \newcommand{\coverBodyFont}{\rmfamily}
}{}

% ======================================================================
% Profile: garamond (trade / pop-sci voice)
% EB Garamond body; Source Sans 3 companion sans with a documented
% fallback chain (Source Sans 3 -> Source Sans Pro -> TeX Gyre Heros),
% because Source Sans is the one stack member frequently absent from
% stock TeX installs; Noto Sans Mono for code snippets in prose.
% ======================================================================
\ifdefstring{\BookFontProfile}{garamond}{%
  \setmainfont{EB Garamond}[Numbers=OldStyle]
  \IfFontExistsTF{Source Sans 3}{%
    \setsansfont{Source Sans 3}%
  }{%
    \IfFontExistsTF{Source Sans Pro}{%
      \setsansfont{Source Sans Pro}%
    }{%
      % TeX Gyre Heros (Helvetica clone) ships with TeX Live; scaled
      % down because its x-height towers over Garamond's.
      \setsansfont{TeX Gyre Heros}[Scale=0.88]%
    }%
  }%
  \IfFontExistsTF{Noto Sans Mono}{%
    % Real bold exists; italic doesn't — fake the slant for comments.
    \setmonofont{Noto Sans Mono}[Scale=MatchLowercase,AutoFakeSlant=0.18]%
  }{%
    \setmonofont{TeX Gyre Cursor}[Scale=MatchLowercase]% Courier clone fallback
  }%
  \AtBeginEnvironment{tabular}{\addfontfeatures{Numbers={Monospaced,Lining}}}
  \newcommand{\coverTitleFont}{\rmfamily\bfseries}
  \newcommand{\coverSubtitleFont}{\sffamily}
  \newcommand{\coverBodyFont}{\rmfamily}
}{}

% ======================================================================
% Profile: plex (technical voice — RFC/hacking/vibe-coding books)
% Optical calibration notes (from the RFC book's print proofs):
%   * Plex true Bold is too heavy on cream 60# stock at 10-11pt —
%     SemiBold is mapped as bold everywhere instead.
%   * Plex has NO small caps (smcp) and NO oldstyle figures (onum);
%     \bkscshape/\bklabelcaps degrade to letterspaced-free uppercase,
%     and figures stay lining (they are lining-by-design, tabular in
%     the mono; no tabular hook needed).
% ======================================================================
\ifdefstring{\BookFontProfile}{plex}{%
  % SemiBold faces are addressed by POSTSCRIPT NAME (IBMPlexSans-SemiBold),
  % not display name: fontconfig indexes Plex SemiBold under the mangled
  % family "IBM Plex Sans SmBld" while texmf uses "... SemiBold", so the
  % display name resolves on some machines and not others. The PS name is
  % embedded in the font file itself and identical in the TTF and OTF
  % releases — still a font NAME, never a file path (ADR 0004).
  \setmainfont{IBM Plex Serif}[
    BoldFont       = {IBMPlexSerif-SemiBold},
    BoldItalicFont = {IBMPlexSerif-SemiBoldItalic},
  ]
  \setsansfont{IBM Plex Sans}[
    BoldFont       = {IBMPlexSans-SemiBold},
    BoldItalicFont = {IBMPlexSans-SemiBoldItalic},
  ]
  \setmonofont{IBM Plex Mono}[
    Scale          = MatchLowercase,
    BoldFont       = {IBMPlexMono-SemiBold},
    BoldItalicFont = {IBMPlexMono-SemiBoldItalic},
  ]
  % No smcp in Plex: small-caps contexts become plain uppercase.
  \renewcommand{\bkscshape}{\upshape}
  \renewcommand{\bklabelcaps}[1]{{\small\sffamily\MakeUppercase{#1}}}
  \newcommand{\coverTitleFont}{\sffamily\bfseries}% Plex Sans SemiBold display
  \newcommand{\coverSubtitleFont}{\sffamily}
  \newcommand{\coverBodyFont}{\rmfamily}
}{}

% ======================================================================
% Per-profile microtype protrusion (hanging punctuation).
% IMPLEMENTATION NOTE: these blocks MUST sit at top level. microtype's
% list parser re-scans its argument, and dies with "Unknown slot
% number" / "Improper alphabetic constant" if the list has already
% been frozen as a macro argument — so no \ifdefstring{...}{...}
% wrappers here; primitive \ifx / \ifdefined only (they skip tokens
% without grabbing them). \ifdefined\SetProtrusion doubles as the
% "is microtype loaded" guard for cover.tex contexts.
% ======================================================================
\ifdefined\SetProtrusion
\makeatletter
% \newcommand (not \def): \ifx also compares the \long prefix, and
% \BookFontProfile comes from \newcommand in generated/metadata.tex.
\newcommand{\bk@fontprofile@libertinus}{libertinus}
\newcommand{\bk@fontprofile@garamond}{garamond}
\newcommand{\bk@fontprofile@plex}{plex}

\ifx\BookFontProfile\bk@fontprofile@libertinus
  % Gentle trade-book edge hang; Libertinus' defaults are already fair.
  \SetProtrusion[load=T1-default]{encoding=TU}{
    . = {0,600}, {,} = {0,600}, : = {0,400}, ; = {0,400},
    - = {0,600}, \textendash = {0,300}, \textemdash = {0,250},
    \textquoteleft = {600,0}, \textquoteright = {0,600},
  }
\fi

\ifx\BookFontProfile\bk@fontprofile@garamond
  % Garamond's delicate serifs reward stronger period/comma hang.
  \SetProtrusion[load=T1-default]{encoding=TU}{
    . = {0,700}, {,} = {0,700}, : = {0,400}, ; = {0,400},
    - = {0,600}, \textendash = {0,300}, \textemdash = {0,250},
    \textquoteleft = {700,0}, \textquoteright = {0,700},
  }
\fi

\ifx\BookFontProfile\bk@fontprofile@plex
  % RFC book's hand-tuned table (incl. dashes, parens, slash — Plex's
  % wide punctuation needs the stronger values).
  \SetProtrusion[load=T1-default]{encoding=TU}{
    . = {0,700}, {,} = {0,700}, : = {0,500}, ; = {0,500},
    - = {0,700}, \textendash = {0,300}, \textemdash = {0,300},
    ( = {400,0}, ) = {0,400}, / = {200,300},
    \textquoteleft = {700,0}, \textquoteright = {0,700},
  }
\fi
\makeatother
\fi

% ----------------------------------------------------------------------
% CJK / emoji fallback hook (LuaLaTeX only — ADR 0004).
% Disabled by default: fallback fonts are machine-dependent and color
% emoji is rarely print-appropriate. To enable, define
% \BookEnableFontFallback in pretex (or here) and adjust the families;
% luaotfload then substitutes glyphs the body face lacks.
% ----------------------------------------------------------------------
\ifluatex
  \ifdefined\BookEnableFontFallback
    \directlua{luaotfload.add_fallback("bookfallback", {
      "Noto Sans CJK SC:mode=node;script=hani;",
      "Noto Color Emoji:mode=harf;",
    })}
    % NOTE: fallbacks attach per-font; re-declare the main font with
    % RawFeature={fallback=bookfallback} in the active profile branch.
  \fi
\fi
